% TOP_PROBLEM: {"id":30007602,"problem_number":"LOCAL-30007602","title":"Optimal capital requirements","category_id":21,"category":{"id":21,"name":"economics","display_name":"Economics","description":"Open economic theory statements, published research agendas, and proposed scoped specifications; question provenance and historical priority are qualified per record.","slug":"economics","order_index":21,"created_at":"2026-10-08T00:00:00Z"},"status":"research_question","external_url":null,"legacy_ids":[],"published":false,"statement":"In the explicitly specified setting: What bank-capital structure and buffer levels maximize social welfare after accounting for crisis prevention, credit supply and regulatory complexity? The mathematical statement above fixes the target, assumptions and scope of this catalogue formulation.","economics":{"original_id":91,"field":"Banking, credit and financial stability","kind":"design","formalization_basis":"adapted_research_question","source_status":"research_agenda","priority_status":"earliest_found","priority_note":"This is the earliest explicit statement verified in this audit; first priority for the full catalogue formulation is not established.","earliest_verified_reference":{"authors":["Bank of England"],"title":"Bank of England Agenda for Research, 2025–2028","year":2026,"url":"https://www.bankofengland.co.uk/research/bank-of-england-agenda-for-research","doi":null,"locator":"Prudential Architecture / Evaluating regulations; Regulatory framework, paragraph 1","evidence_summary":"Questions regulatory cost-benefit analysis and complex capital-stack interactions."},"current_reference":{"authors":["Bank of England"],"title":"Bank of England Agenda for Research, 2025–2028","year":2026,"url":"https://www.bankofengland.co.uk/research/bank-of-england-agenda-for-research","doi":null,"locator":"Prudential Architecture / Evaluating regulations; Regulatory framework, paragraph 1","evidence_summary":"Questions regulatory cost-benefit analysis and complex capital-stack interactions."},"formalization_note":"Adapts an explicitly verified research-agenda passage. Mathematical objects, interventions, objectives and scope are proposed here; existing model-specific solutions are not relabeled unsolved."},"statusQualification":"Published question or research agenda verified at the cited locator. The mathematical specification is adapted for this catalogue and includes additional assumptions; source publication does not establish first-ever priority or certify this exact model remains unsolved. Adapts an explicitly verified research-agenda passage. Mathematical objects, interventions, objectives and scope are proposed here; existing model-specific solutions are not relabeled unsolved.","statusReviewedAt":"2026-10-08"}
% FIRST_POSTED: 2026-10-08
% ENTRY_KIND: statement_only
% !TeX program = lualatex
\documentclass[11pt]{article}
\usepackage{fontspec}
\usepackage[a4paper,margin=25mm]{geometry}
\usepackage{amsmath,amssymb,mathtools}
\usepackage{xurl}
\usepackage[hidelinks]{hyperref}
\newcommand{\sref}[2]{\hyperlink{source:#1:#2}{[S#2]}}
\setlength{\emergencystretch}{3em}
\begin{document}
% problemId:30007602
\section*{Optimal capital requirements}\label{30007602}
\subsection{Definitions and mathematical statement}
Consider banks that finance loans and securities with deposits, debt and equity. A regulatory vector \(k\in K\) specifies risk-weighted capital, leverage, stress buffers and release rules. Households and banks choose consumption, lending and financing in a dynamic equilibrium with limited liability and explicitly stated public guarantees. Let \(W(k;P)\) be discounted social welfare under primitive shock law \(P\), including borrower surplus, household consumption, fiscal costs and default externalities; equity transfers are not automatically resource costs. Let \(\mathcal P\) be a documented set of admissible shock laws and \(K\) a feasible compact set of rules. The task is to characterize
\[k^*\in\arg\max_{k\in K}\inf_{P\in\mathcal P}W(k;P).\]
Quantify how crisis-frequency reductions compare with ordinary lending and intermediation costs, and identify the marginal welfare effects of each component of the regulatory stack. Include endogenous risk selection, bank entry and displacement of lending outside banks. Identify or bound welfare derivatives using independently disciplined preference and technology parameters, with explicit uncertainty about rare crises. State which constraints are binding under the proposed optimum. An optimum in a representative-bank model does not establish a universal capital requirement for heterogeneous systems.

\par\textbf{Formulation and scope.} Adapts an explicitly verified research-agenda passage. Mathematical objects, interventions, objectives and scope are proposed here; existing model-specific solutions are not relabeled unsolved.

\par\textbf{Open-question provenance.} An explicit question or research agenda was verified at the source locator. This does not by itself certify that every later version remains unresolved \sref{30007602}{1}.

\par\textbf{Historical priority.} This is the earliest explicit statement verified in this audit; first priority for the full catalogue formulation is not established.

\subsection{Short English statement}
In the explicitly specified setting: What bank-capital structure and buffer levels maximize social welfare after accounting for crisis prevention, credit supply and regulatory complexity? The mathematical statement above fixes the target, assumptions and scope of this catalogue formulation.

\subsection{Sources}
\begin{itemize}\raggedright
\item[\textnormal{[S1]}]\hypertarget{source:30007602:1}{} Bank of England. 2026. Bank of England Agenda for Research, 2025–2028. Prudential Architecture / Evaluating regulations; Regulatory framework, paragraph 1.
\url{https://www.bankofengland.co.uk/research/bank-of-england-agenda-for-research}.
{\small\textit{Earliest verified question/agenda reference; historical priority qualified below. Questions regulatory cost-benefit analysis and complex capital-stack interactions.}}
\end{itemize}
\end{document}
